% ============================================================================
%  บทที่ 4: การสร้างและบูรณาการระบบ
% ============================================================================

\chapter{การสร้างและบูรณาการระบบ}
\label{ch:implementation}

% ============================================================================
\section{การประกอบโครงสร้างทางกล}
\label{sec:mech_assembly}
% ============================================================================

\subsection{ภาพรวมการประกอบ}
\label{subsec:mech_overview}

การประกอบระบบเชิงกลแบ่งออกเป็น 3 Sub-assembly หลัก
ได้แก่ ส่วนฐาน (Base), ส่วนแขนกลและข้อต่อ (Arm Links \& Joints)
และชุดเอ็นโค้ดเดอร์ (Encoder Sub-assembly)
ก่อนประกอบรวมเป็นหุ่นยนต์เต็มรูปแบบ

% -----------------------------------------------------------------------
% รูปที่: Exploded View (Isometric)
% ชื่อไฟล์แนะนำ: fig_exploded_iso.png
% -----------------------------------------------------------------------
% \begin{figure}[H]
%     \centering
%     \includegraphics[width=0.88\textwidth]{fig_exploded_iso}
%     \caption{ภาพรวมชิ้นส่วนของหุ่นยนต์ (Exploded View) มุม Isometric}
%     \label{fig:exploded_iso}
% \end{figure}

% -----------------------------------------------------------------------
% รูปที่: Exploded View (Side)
% ชื่อไฟล์แนะนำ: fig_exploded_side.png
% -----------------------------------------------------------------------
% \begin{figure}[H]
%     \centering
%     \includegraphics[width=0.88\textwidth]{fig_exploded_side}
%     \caption{ภาพรวมชิ้นส่วนของหุ่นยนต์ (Exploded View) มุม Side View}
%     \label{fig:exploded_side}
% \end{figure}

\subsection{การประกอบ Sub-assembly}
\label{subsec:sub_assembly}

\subsubsection{ส่วนฐาน (Base Sub-assembly)}

\begin{table}[H]
\centering
\caption{รายการชิ้นส่วนสำหรับประกอบส่วนฐาน}
\label{tab:base_parts}
\begin{tabular}{llc}
\toprule
ชิ้นส่วน & รายละเอียด & จำนวน \\
\midrule
อลูมิเนียมโปรไฟล์ 20$\times$40 & ยาว 100\,mm & 4 \\
แผ่นฐานหุ่น (Base Plate)    & AISI 316 & 1 \\
แผ่นรอง Pulley              & AISI 316 & 1 \\
Pulley                       & 70 ฟัน, ID 19\,mm & 1 \\
Bearing Mounting             & เหล็ก & 1 \\
Tapered Roller Bearing       & ID 17\,mm, OD 40\,mm & 1 \\
ชิ้นส่วน Set Home            & เหล็ก & 1 \\
Ball Bearing                 & ID 12\,mm, OD 19\,mm & 1 \\
Socket Head Cap Screw M5     & ยาว 12\,mm & 16 \\
Socket Head Cap Screw M4     & ยาว 40\,mm & 4 \\
Nut M4                       & -- & 4 \\
\bottomrule
\end{tabular}
\end{table}

ขั้นตอนการประกอบส่วนฐาน:
\begin{enumerate}
    \item ประกอบอลูมิเนียมโปรไฟล์ 4 ชิ้นเข้ากับ Base Plate
          ด้วยสกรู M5 ยาว 12\,mm จำนวน 8 ตัว
    \item ประกอบแผ่นรอง Pulley ปิดทับด้านบน
          ด้วยสกรู M5 ยาว 12\,mm จำนวน 8 ตัว
    \item วาง Pulley 70 ฟัน ลงในตำแหน่งแกนกลาง
    \item ยึด Bearing Mounting ด้วยสกรู M4 ยาว 40\,mm
          จำนวน 3 ตัว พร้อม Nut
    \item ใส่ Tapered Roller Bearing (ด้านบน)
          และ Ball Bearing (ด้านล่าง)
    \item ติดตั้งชิ้นส่วน Set Home
          ด้วยสกรู M4 ยาว 40\,mm พร้อม Lock Nut
\end{enumerate}

\subsubsection{ส่วนแขนกลและข้อต่อ (Arm Sub-assembly)}

\begin{table}[H]
\centering
\caption{รายการชิ้นส่วนสำหรับประกอบส่วนแขนกล}
\label{tab:arm_parts}
\begin{tabular}{llc}
\toprule
ชิ้นส่วน & รายละเอียด & จำนวน \\
\midrule
Motor                        & RS-775 DC Motor & 1 \\
แผ่นแขนกล (บน)              & Aluminum Alloy 6061 & 1 \\
แผ่นแขนกล (ล่าง)            & Aluminum Alloy 6061 & 1 \\
เหล็กกล่อง Galvanized       & 25$\times$25\,mm & 1 \\
Stepped Shaft                & Stainless Steel & 1 \\
Pulley                       & 24 ฟัน, ID 12\,mm & 1 \\
Idler                        & ID 10\,mm, Width 15\,mm & 1 \\
Shaft Support                & Aluminum Alloy 10\,mm & 3 \\
Key (ลิ่ม)                  & 5$\times$5$\times$8\,mm & 1 \\
Proximity Sensor             & $\varnothing$12\,mm & 1 \\
\bottomrule
\end{tabular}
\end{table}

ขั้นตอนการประกอบส่วนแขนกล:
\begin{enumerate}
    \item ประกบแผ่นแขนกลบน-ล่างเข้ากับเหล็กกล่อง Galvanized
          พร้อมบูชรองรักษาระยะ
    \item ติดตั้งมอเตอร์และสวม Pulley 24 ฟัน เข้ากับเพลามอเตอร์
    \item สอด Stepped Shaft พร้อม Key ที่จุดหมุนหลัก
          และติดตั้ง Shaft Support
    \item ติดตั้ง Idler และ Proximity Sensor
    \item ขันน็อตทุกจุดและตรวจสอบการเคลื่อนที่
\end{enumerate}

\subsubsection{ชุดเอ็นโค้ดเดอร์ (Encoder Sub-assembly)}

\begin{table}[H]
\centering
\caption{รายการชิ้นส่วนสำหรับประกอบชุดเอ็นโค้ดเดอร์}
\label{tab:encoder_parts}
\begin{tabular}{llc}
\toprule
ชิ้นส่วน & รายละเอียด & จำนวน \\
\midrule
AMT Mounting          & AISI 316 & 1 \\
Encoder AMT-103       & 2048 PPR & 1 \\
ชิ้นส่วนประกบ (บน/ล่าง) & 3D Print & อย่างละ 1 \\
บูชรอง M3             & OD 5.5$\times$20\,mm & 4 \\
สกรู M2/M3/M5         & ตามกำหนด & 16 \\
\bottomrule
\end{tabular}
\end{table}

\subsection{การประกอบรวมขั้นสุดท้าย}
\label{subsec:final_assembly}

\begin{enumerate}
    \item ติดตั้งชุดเอ็นโค้ดเดอร์เข้ากับฐาน
    \item คล้องสายพาน HTD5M เตรียมไว้
    \item สอดเพลาแขนกลลงในรูร้อยแกนกลางของฐาน
    \item ตั้งค่า Bearing Preload โดยขันน็อตเพลา
          เพื่อกดทับ Tapered Roller Bearing (ลดระยะคลอน)
    \item ต่อเพลาเชื่อมต่อ Field Encoder ด้วย Shaft Coupling
    \item ติดตั้ง Gripper ที่ปลายแขน ยึดด้วย Lock Nut 3 ตัว
\end{enumerate}

\subsection{ปัญหาที่พบระหว่างการประกอบ}
\label{subsec:assembly_issues}

\begin{itemize}
  \item \textbf{รู Proximity Sensor ขนาดไม่ตรง:}
        รูออกแบบ $\varnothing$8\,mm แต่ Sensor จริง $\varnothing$12\,mm
        แก้ไขโดยเจาะขยายรูเพิ่มใน Lab

  \item \textbf{รู Driver Pulley และ Arm เบี้ยว:}
        ชิ้นส่วนที่ส่งผลิตภายนอกมีรูไม่ตรงแนวศูนย์
        เนื่องจาก Setup งานไม่ตรงกับแบบ
        แก้ไขโดยนำกลับไปปรับแต่งเพิ่มเติม

  \item \textbf{แผ่นรอง Pulley ไม่ได้ปาดผิว:}
        ผู้ผลิตไม่ได้ปาดผิวและร่อง Aluminium Profile
        ตามแบบที่กำหนด ดำเนินการแก้ไขเองใน Lab
\end{itemize}

% -----------------------------------------------------------------------
% รูปที่: CAD Assembly
% ชื่อไฟล์แนะนำ: fig_cad_assembly.png
% -----------------------------------------------------------------------
% \begin{figure}[H]
%     \centering
%     \includegraphics[width=0.85\textwidth]{fig_cad_assembly}
%     \caption{CAD Assembly ของระบบเชิงกลแสดงชิ้นส่วนหลักทั้งหมด}
%     \label{fig:cad_assembly}
% \end{figure}

% ============================================================================
\section{การประกอบตู้ควบคุมและการเดินสายไฟ}
\label{sec:elec_assembly}
% ============================================================================

\subsection{Bill of Materials ของตู้ควบคุม}
\label{subsec:bom}

\begin{table}[H]
  \centering
  \caption{Bill of Materials (BOM) ของตู้ควบคุม}
  \label{tab:bom}
  \begin{tabular}{llcc}
    \toprule
    Reference & อุปกรณ์ & จำนวน & หมายเหตุ \\
    \midrule
    U1A  & STM32G474RE Nucleo Board  & 1 & Main Controller \\
    U2   & ESP32-DEVKIT-V1            & 1 & Joystick Receiver \\
    --   & Cytron MD10C Motor Driver  & 1 & 24V/10A \\
    CB1  & Circuit Breaker 10A        & 1 & AC Input \\
    CB2  & Circuit Breaker 5A         & 1 & 24V Bus \\
    CB3  & Circuit Breaker 3A         & 1 & 5V Bus \\
    --   & PSU 24V/10A                & 1 & Motor Power \\
    --   & PSU 5V/5A                  & 1 & Logic Power \\
    --   & Magnetic Contactor         & 1 & E-Stop Interlock \\
    --   & Optocoupler Module 8-ch    & 1 & Signal Isolation \\
    --   & Relay Module 4-ch          & 1 & Output Control \\
    --   & Emergency Stop Button      & 1 & NC Contact \\
    --   & Pilot Lamp Blue (Power)    & 1 & Power Status \\
    --   & Pilot Lamp Red (Emergency) & 1 & Fault Status \\
    --   & Pilot Lamp Green (Ready)   & 1 & Ready Status \\
    --   & Reset Button               & 1 & NC Contact \\
    --   & Mode Selector Switch       & 1 & Manual/Auto \\
    --   & Fuse (Motor Circuit)       & 1 & Short Circuit \\
    \bottomrule
  \end{tabular}
\end{table}

\subsection{แนวทางการเดินสายไฟ}
\label{subsec:cable_management}

การเดินสายแบ่งออกเป็น 2 วงจรหลักที่แยก Physical กันอย่างชัดเจน
สัญญาณดิจิทัลทุกเส้นระหว่าง High Voltage Side และ Logic Side
ผ่าน Optocoupler Module เพื่อแยก Ground และป้องกัน Noise

\begin{itemize}
  \item \textbf{Service Loop:} บริเวณจุดหมุนของแขนกล
        เผื่อความยาวสายเพื่อให้โค้งงอได้อิสระ
        ป้องกันสายดึงรั้งหรือขาดเมื่อหมุนถึงมุมสูงสุด

  \item \textbf{Bending Radius:} ติดตั้งชิ้นส่วน 3D Print
        ที่มุมของหุ่นยนต์เพื่อบังคับให้สายโค้งงอในรัศมีที่เพียงพอ
        ป้องกันการหักมุมแคบที่ทำลายสายสัญญาณภายใน
\end{itemize}

% -----------------------------------------------------------------------
% รูปที่: ผลการเดินสายภายในตู้ควบคุม
% ชื่อไฟล์แนะนำ: fig_cabinet_assembled.png
% -----------------------------------------------------------------------
% \begin{figure}[H]
%     \centering
%     \includegraphics[width=0.78\textwidth]{fig_cabinet_assembled}
%     \caption{ผลการเดินสายภายในตู้ควบคุมไฟฟ้า}
%     \label{fig:cabinet_assembled}
% \end{figure}

% ============================================================================
\section{การพัฒนา Firmware}
\label{sec:firmware}
% ============================================================================

\subsection{โครงสร้างโมดูลและ Task Timing}
\label{subsec:firmware_modules}

Firmware พัฒนาบน STM32CubeIDE ใช้ HAL Library
แบ่งเป็น 4 โมดูลหลัก ดังแสดงใน\cref{tab:firmware_modules}

\begin{table}[H]
  \centering
  \caption{โมดูลหลักใน STM32 Firmware}
  \label{tab:firmware_modules}
  \begin{tabular}{lp{6.5cm}l}
    \toprule
    โมดูล & หน้าที่ & ไฟล์ \\
    \midrule
    \texttt{motor\_controller} &
      Encoder อ่านตำแหน่ง/ความเร็ว,
      Cascade PID (Position + Velocity),
      S-Curve Trajectory Generator,
      Kalman Filter (Steady-State),
      ZVD Input Shaper (3-Impulse, Circular Buffer),
      Joystick Command Parser,
      Safety Monitor (Stall, E-Stop) &
      \texttt{motor\_controller.c} \\
    \addlinespace
    \texttt{modbus\_bridge} &
      Modbus Register Mapping (128 registers),
      Command Dispatch จาก Master,
      Status Update ไปยัง Register Frame &
      \texttt{modbus\_bridge.c} \\
    \addlinespace
    \texttt{modbus\_rtu} &
      Modbus RTU Protocol Engine,
      CRC16 Verification,
      Frame Parser (FC03/FC06),
      FSM (IDLE/RECEPTION/PROCESSING/EMISSION) &
      \texttt{modbus\_rtu.c} \\
    \addlinespace
    \texttt{main} &
      Peripheral Init, Callback Registration,
      Main Loop (Modbus Process,
      Register Update, MATLAB Data Send) &
      \texttt{main.c} \\
    \bottomrule
  \end{tabular}
\end{table}

\subsubsection{Execution Model และ Task Timing}
\label{subsubsec:execution_model}

TIM6 ตั้งค่าที่ \textbf{1\,kHz} (Prescaler = 169, Period = 169,
ที่ System Clock 170\,MHz ได้ Period จริง $(169+1)(169+1)/170\,\text{MHz} = 1$\,ms)
การแบ่ง Loop สองความถี่ทำผ่าน Software Counter ภายใน ISR เดียว:

\begin{itemize}
  \item \textbf{ทุก Tick (1\,kHz):}
        \texttt{MotorCtrl\_Tick1kHz()} รัน
        อ่าน TIM3 Encoder $\rightarrow$ Kalman Predict+Update $\rightarrow$
        Velocity Inner PID $\rightarrow$ \texttt{Motor\_SetPWM()}
  \item \textbf{ทุก 10th Tick (100\,Hz):}
        \texttt{HwIo\_Poll100Hz()} debounce สัญญาณ I/O\\
        \texttt{MotorCtrl\_Tick100Hz()} รัน
        S-Curve Integrator $\rightarrow$ ZVD Shaper $\rightarrow$
        Position Outer PID $\rightarrow$ Velocity Setpoint สำหรับ Inner Loop
\end{itemize}

สถาปัตยกรรมนี้ทำให้ Velocity Loop สามารถรัน 10 ครั้งต่อ 1 Position Loop cycle
ซึ่งรับประกัน Bandwidth Separation ที่เพียงพอโดยไม่ต้องใช้ Timer แยก

\begin{table}[H]
  \centering
  \caption{ISR และ Task Timing ของ Firmware}
  \label{tab:isr_timing}
  \small
  \begin{tabular}{llll}
    \toprule
    Timer/Peripheral & Period & หน้าที่ & Callback \\
    \midrule
    TIM6  & \textbf{1\,ms (1\,kHz)} &
      Control Loop Base Rate &
      \texttt{Timer\_PeriodElapsedCallback} \\
    & \quad \textit{ทุก tick} & Velocity Inner PID + Kalman & \textit{(ใน TIM6 ISR)} \\
    & \quad \textit{ทุก 10th tick} & Position Outer PID + ZVD & \textit{(ใน TIM6 ISR)} \\
    TIM16 & 3\,ms &
      Modbus T3.5 Timeout &
      \texttt{Timer\_PeriodElapsedCallback} \\
    TIM3  & -- (QEI) &
      Quadrature Encoder Count &
      Polling ใน TIM6 ISR \\
    TIM1  & $f_{PWM} \approx 19.6$\,kHz &
      Motor PWM Output (ARR=50) &
      -- \\
    LPUART1 & -- (UART RX IT) &
      Modbus RTU (230400, 8N1) &
      \texttt{Modbus\_RxCallback} \\
    USART3  & -- (UART RX IT) &
      Joystick Command (115200, 8N1) &
      \texttt{Joystick\_RxCallback} \\
    \bottomrule
  \end{tabular}
\end{table}

\subsubsection{PID Gains ที่ใช้ใน Firmware จริง}
\label{subsubsec:pid_gains_firmware}

ค่า PID Gain ที่ระบุในตารางต่อไปนี้คือค่าที่วัดจาก Hardware โดยตรง
(\texttt{params.h}) ซึ่งต่างจากค่าที่ได้จาก Root Locus Design (บทที่~3)
เนื่องจาก Reduced-Order Model มีความแตกต่างจาก Plant จริง
และต้องผ่านกระบวนการ Hardware Commissioning

\begin{table}[H]
  \centering
  \caption{PID Gains ที่ใช้ใน Firmware จริง (Hardware-Identified)}
  \label{tab:pid_gains_hw}
  \begin{tabular}{llccc}
    \toprule
    Loop & Gain & สัญลักษณ์ & ค่า (Hardware) & หน่วย \\
    \midrule
    \multirow{3}{*}{Velocity Inner}
      & Proportional & $K_{p,vel}$ & 1.58  & PWM/(rad/s) \\
      & Integral     & $K_{i,vel}$ & 0.50  & PWM/rad \\
      & Derivative   & $K_{d,vel}$ & 0.00  & -- \\
    \midrule
    \multirow{3}{*}{Position Outer}
      & Proportional & $K_{p,pos}$ & 3.00  & (rad/s)/rad \\
      & Integral     & $K_{i,pos}$ & 0.02  & -- \\
      & Derivative   & $K_{d,pos}$ & 0.15  & -- \\
    \midrule
    \multirow{2}{*}{Feedforward}
      & Velocity FF  & $K_{vff}$   & 3.03  & PWM/(rad/s) \\
      & Accel FF     & $K_{aff}$   & 0.10  & PWM/(rad/s²) \\
    \bottomrule
    \multicolumn{5}{l}{\small หมายเหตุ: ค่าเหล่านี้วัดจาก Hardware โดยตรง} \\
    \multicolumn{5}{l}{\small ห้ามเปลี่ยนโดยไม่ผ่านกระบวนการ System Identification ใหม่} \\
  \end{tabular}
\end{table}

\subsection{การทำงานร่วมกันของโมดูล}
\label{subsec:module_interaction}

การทำงานแบ่งออกเป็น 2 Context ดังนี้

\noindent\textbf{Interrupt Context} (Priority สูง):
\texttt{Modbus\_RxCallback} รับ Byte ทีละตัวจาก LPUART1,
\texttt{Joystick\_RxCallback} รับ Command จาก USART3
และเรียก \texttt{Motor\_ProcessCommand} ทันที
และ TIM6 เรียก \texttt{Motor\_ControlLoop} ทุก 10\,ms

\medskip
\noindent\textbf{Main Loop Context} (Priority ต่ำ):
\texttt{ModbusBridge\_UpdateRegisters} ซิงค์ตำแหน่งและความเร็วลง Register Frame,
\texttt{ModbusBridge\_Process} รัน Protocol Engine
และ \texttt{Motor\_SendDataToMatlab} ส่งข้อมูลออก USART3
ทุก 20\,ms สำหรับ Live Monitoring

\subsection{การ Implement ZV Input Shaper (Circular Buffer)}
\label{subsec:zv_impl}

ZV Input Shaper ถูก Implement บน STM32 ด้วย Circular Buffer
โดย Setpoint $r(t)$ ถูก Delay เป็นเวลา $t_2 = 0.259$\,s
(= 26 Control Cycles ที่ 100\,Hz) แล้วรวมกับ

\begin{equation}
  r_{shaped}(t) = A_1 \cdot r(t) + A_2 \cdot r(t - t_2)
  \label{eq:zv_impl}
\end{equation}

Buffer มีขนาด $\lceil t_2 / T_p \rceil + 1 = 27$ Elements
ทำงานที่ Position Outer Loop (100\,Hz) ใน \texttt{Motor\_ControlLoop}

\subsection{การ Implement Kalman Filter (Steady-State)}
\label{subsec:kf_impl}

Kalman Filter 4 สถานะ $\mathbf{x} = [\theta, \omega, i_a, \tau_d]^\top$
ใช้ Steady-State Gain $\mathbf{K}_\infty$ ที่คำนวณล่วงหน้าด้วย DARE บน MATLAB
แล้ว Hardcode ใน Firmware เพื่อลด Computational Load

ทุก Control Cycle (10\,ms) ทำงาน 2 ขั้นตอน:

\begin{enumerate}
  \item \textbf{Predict:}
        $\hat{\mathbf{x}}_{k|k-1} = \mathbf{A}_d\,\hat{\mathbf{x}}_{k-1} + \mathbf{B}_d\,u_{k-1}$
  \item \textbf{Update:}
        $\hat{\mathbf{x}}_{k|k} = \hat{\mathbf{x}}_{k|k-1} + \mathbf{K}_\infty(y_k - \mathbf{C}\hat{\mathbf{x}}_{k|k-1})$
\end{enumerate}

Output $\hat{\theta}$ และ $\hat{\omega}$ ใช้ใน Control Loop
ส่วน $\hat{\tau}_d$ ใช้สำหรับ Monitoring

\subsection{การ Implement Cascade PID (Discrete)}
\label{subsec:pid_impl}

PID Discretize ด้วยวิธี Tustin (Bilinear) ตาม\cref{eq:tustin}

Velocity Inner Loop (1\,kHz) ใช้ PI + Anti-windup Back-calculation

\begin{equation}
  u_v[k] = K_{p,v}\,e_v[k] + \text{integ}[k] + K_{d,v}\frac{e_v[k]-e_v[k-1]}{T_v}
  \label{eq:pid_discrete_v}
\end{equation}

\begin{equation}
  \text{integ}[k] = \text{integ}[k-1] + K_{i,v}\,e_v[k]\,T_v
                  + K_{t,aw}(u_{sat}[k] - u_{pre}[k])\,T_v
  \label{eq:anti_windup_discrete}
\end{equation}

% ============================================================================
\section{การบูรณาการระดับระบบ}
\label{sec:integration}
% ============================================================================

\subsection{การเชื่อมต่อ Motor, Encoder และ Motor Driver}
\label{subsec:motor_integration}

Encoder AMT-103 (2048 PPR $\times$ 4 = 8192 counts/rev)
ต่อเข้า STM32 TIM3 ผ่านโหมด Quadrature Encoder Interface (QEI)
บน Pin PB4 (CH A) และ PB5 (CH B)

Motor Driver Cytron MD10C รับสัญญาณ PWM จาก TIM8 CH2 (PA8)
และสัญญาณ DIR จาก PA9 ผ่าน Optocoupler เพื่อแยก Ground

\subsection{การเชื่อมต่อ Joystick ผ่าน ESP32}
\label{subsec:joystick_integration}

ESP32-DEVKIT-V1 รับสัญญาณ Bluetooth จาก Joystick
แล้วส่ง Command Character ผ่าน USART3 (115200 bps, 8N1)
ไปยัง STM32 ที่ Pin PC10 (TX) และ PC11 (RX)

\begin{table}[H]
  \centering
  \caption{การ Map Command Character กับคำสั่งระบบ}
  \label{tab:joystick_map_fw}
  \begin{tabular}{llll}
    \toprule
    Character & คำสั่ง & โหมดที่ใช้ได้ & หมายเหตุ \\
    \midrule
    \texttt{L} & หมุนซ้าย (Step/Jog)  & Manual & Coarse/Fine \\
    \texttt{R} & หมุนขวา (Step/Jog)   & Manual & Coarse/Fine \\
    \texttt{O} & Release / Stop Jog   & Manual & -- \\
    \texttt{A} & Home (Hold 3\,s)     & ทุกโหมด & Hold $<$3\,s = Reset Encoder \\
    \texttt{P} & E-Stop / Reset       & ทุกโหมด & Toggle E-Stop \\
    \texttt{M} & สลับ Coarse/Fine     & Manual & Toggle Jog Mode \\
    \texttt{U} & Gripper Up           & Manual & -- \\
    \texttt{D} & Gripper Down         & Manual & -- \\
    \texttt{F} & Gripper Toggle       & Manual & Open/Close \\
    \texttt{S} & Gripper Pick         & Manual & -- \\
    \texttt{G} & Gripper Place        & Manual & -- \\
    \texttt{J} & Hold Position        & Manual & Lock ตำแหน่งปัจจุบัน \\
    \bottomrule
  \end{tabular}
\end{table}

การสลับโหมดระหว่าง Joystick และ Base System ทำผ่าน B-Button Hold 1\,s
ซึ่งตั้งค่า \texttt{control\_system\_mode}
เมื่ออยู่ใน \texttt{CONTROL\_MODE\_BASE\_SYSTEM}
คำสั่งจาก Joystick ทุกตัวถูก Block ยกเว้น
\texttt{P} (E-Stop), \texttt{A} (Home) และ \texttt{M} (Mode Toggle)

\subsection{การเชื่อมต่อ Modbus RTU และ Base System}
\label{subsec:modbus_integration}

ระบบสื่อสารกับ Base System ผ่าน Modbus RTU
บน LPUART1 (19200 bps, 8E1, RS-485, Slave Address~21)

Base System ใช้ Python asyncio WebSocket Backend
(แทน \texttt{main.exe} เดิม เนื่องจากปัญหา Abnormal Heartbeat ใน Demo~1)
โดย Backend รับ WebSocket Command จาก Frontend
แล้วแปลงเป็น Modbus RTU Frame ส่งผ่าน Serial Port

Modbus Register Map แบ่งเป็น 128 Registers ประกอบด้วย
Status Registers (ตำแหน่ง, ความเร็ว, สถานะ FSM)
และ Command Registers (Target Position, Mode, Gripper Command)